\documentclass[10pt,a4paper]{article}

\usepackage{amsmath,amssymb,amsthm}
\usepackage{mathtools}
\usepackage[utf8]{inputenc}
\usepackage{geometry}
\usepackage{microtype}
\usepackage{hyperref}
\usepackage{enumitem}
\usepackage{booktabs}
\usepackage{xcolor}

\geometry{
  top=1.25in, bottom=1.25in,
  left=1.35in, right=1.35in
}

\hypersetup{colorlinks=true,linkcolor=black,citecolor=black,urlcolor=black}

\theoremstyle{plain}
\newtheorem{theorem}{Theorem}
\newtheorem{proposition}{Proposition}
\newtheorem{lemma}{Lemma}
\newtheorem{corollary}{Corollary}
\theoremstyle{definition}
\newtheorem{definition}{Definition}
\theoremstyle{remark}
\newtheorem{remark}{Remark}

\newcommand{\Lto}{\xrightarrow{\ \mathsf{L}\ }}
\newcommand{\home}{\lfloor}
\newcommand{\Rho}{\rho}
\newcommand{\Pii}{\pi}
\newcommand{\Mu}{\mu}
\newcommand{\Mubar}{\bar{\mu}}
\newcommand{\Rres}{\mathcal{R}}
\newcommand{\Id}{\mathrm{Id}}

\title{\textbf{Carried Coupling}\\[0.4em]
\large Why an Observer Invariant Survives Admissible Collapse}
\author{Flyxion\\
  \small Independent Researcher}
\date{October 2026}

\begin{document}
\maketitle

\begin{abstract}
\noindent
A chain of admissible collapses is written
$\emptyset \Rres \Lto R_1 \Lto R_2 \Lto \cdots \Lto R_n \Lto \home$.
Along it three fibres of a run behave oppositely: the extensional state
$\Rho$ strictly coarsens and drains to $0$; the provenance $\Pii$ propagates
and is never lost; and the coupling $\Mu$ between a result and an observer is
invariant. We isolate the mechanism that lets $\Mu$ reach the terminal object
$\home$: $\Mu$ factors through $\Pii$ and not through $\Rho$, and $\Pii$ is
carried across each arrow. The result is a short lemma, the
\emph{carried coupling}, seated against the cross-references of
\emph{The Residue of Proof}. The point is not that collapse is harmless but
that it relocates the residue from the lost fibre into the carried one, which
is why a nonzero coupling is preserved to the limit.
\end{abstract}

\section{Setting}

Fix a chain of representations in the sense of
\emph{The Residue of Proof} \cite{residue}, with a declared task family at each
stage. The chain
\[
\emptyset \Rres \Lto R_1 \Lto R_2 \Lto \cdots \Lto R_n \Lto \home
\]
is a sequence of admissible collapses. Each arrow $\Lto$ is a representation
map that is sufficient for its declared family and refusal-respecting; the
terminal $\home$ is the fixed point at which no further collapse is demanded.

Three quantities travel along the chain. Their behaviour is fixed by three
results of the book, which we restate in summary form and then use.

\begin{itemize}[leftmargin=1.4em,nosep]
\item The \emph{extensional state} $\Rho$ is the triple $(D,R,P)$ of derived,
refused, and popped items (Definition 13.3 of \cite{residue},
\emph{Events, histories, state}). It strictly coarsens and drains to $0$: it is
the lossy fibre.
\item The \emph{provenance} $\Pii$ is the set of $\mathsf{Bind}$ events recorded
in the history. It is append-only. It is not a function of the extensional state
(Proposition 13.11 of \cite{residue}, \emph{Proofs are not a function of the
extensional state}).
\item The \emph{coupling} $\Mu(R, \text{you})$ is the quantity of interest
between a representation $R$ and the observer. It is invariant on the chain.
\end{itemize}

\begin{table}[h]
\centering
\begin{tabular}{@{}llll@{}}
\toprule
fibre & object & behaviour & citation \\
\midrule
$\Rho$ & extensional state & strictly coarsens, drains to $0$ & Def.~13.3, Thm.~13.8 \\
$\Pii$ & provenance ($\mathsf{Bind}$ history) & propagates, append-only, never lost & Prop.~13.11 \\
$\Mu$ & observer coupling & invariant & Prop.~1.12, Thm.~10.12 \\
\bottomrule
\end{tabular}
\caption{The three fibres and their opposing behaviour along the $\mathsf{L}$-chain.
Cross-references are to \emph{The Residue of Proof} \cite{residue}.}
\end{table}

\section{Why $\Mu$ factors through $\Pii$ and not through $\Rho$}

Chapter 1 fixes when a task can be read off a representation: a task factors
through a map exactly when it is constant on that map's fibres.

\begin{proposition}[Factorization criterion; Prop.~1.12 of \cite{residue}]
\label{prop:factor}
Let $\varrho : X \to D$ be a representation and $\mathsf{T} : X \to Y$ a task.
There is $\mathsf{T}' : D \to Y$ with $\mathsf{T} = \mathsf{T}' \circ \varrho$
if and only if $\mathsf{T}$ is constant on every fibre of $\varrho$.
\end{proposition}

The extensional state $\Rho$ is sufficient only for tasks constant on its
fibres. The theorem of Chapter 10 makes the converse unconditional: a
non-injective representation is never sufficient for the family of all
$\{0,1\}$-valued tasks.

\begin{theorem}[No free sufficiency; Thm.~10.12 of \cite{residue}]
\label{thm:nofree}
Let $\sigma : H \to R$. Then $\sigma$ is injective if and only if it is
sufficient for the family of all functions $H \to \{0,1\}$. In particular, if
$\sigma$ identifies two distinct histories then some $\{0,1\}$-valued task is
not served.
\end{theorem}

The coupling $\Mu$ distinguishes runs that share an output, so it does not
factor through $\Rho$. This is exactly the failure recorded for the extensional
state in Chapter 13.

\begin{proposition}[Proofs are not a function of the extensional state;
Prop.~13.11 of \cite{residue}]
\label{prop:prov}
There are well-formed histories $H_1, H_2$ with $\Rho(H_1) = \Rho(H_2)$ and
$\Pii(H_1) \ne \Pii(H_2)$, in which the recorded justifications of an item
yield derivations with different leaves. Hence the task that returns a
derivation is not $\Rho$-sufficient, while the extensional state is minimal for
all future behaviour of the search (Thm.~13.8, \emph{Future equivalence}).
\end{proposition}

So $\Mu$ is not constant on the fibres of $\Rho$, and by
Proposition~\ref{prop:factor} it admits no reading from $\Rho$. The proof lives
in $\Pii$, not in the extensional state; the coupling rides the same carried
fibre.

\section{The carried coupling}

Let the chain of collapses act on representations as
$\Rho_0 \supseteq \Rho_1 \supseteq \cdots \supseteq \Rho_n \downarrow 0$ for the
extensional state, and as $\Pii_0 \rightsquigarrow \Pii_1 \rightsquigarrow
\cdots \rightsquigarrow \Pii_n \rightsquigarrow \Pii_\home$ for the provenance,
where each $\rightsquigarrow$ is an append: $\Pii_i \subseteq \Pii_{i+1}$.

We say $\Mu$ \emph{factors through $\Pii$} if there is $\Mubar$ with
$\Mu(R, \text{you}) = \Mubar(\Pii(R), \text{you})$, and that $\Mu$ is
\emph{monotone-carried} if $\Mubar$ does not vanish under the appends that
extend $\Pii$ to $\Pii_\home$: whenever $\Mubar(\Pii_i, \text{you}) \ne 0$ and
$\Pii_i \subseteq \Pii_{i+1}$ then $\Mubar(\Pii_{i+1}, \text{you}) \ne 0$.

\begin{lemma}[Carried coupling]
\label{lem:carried}
Suppose $\Mu$ factors through $\Pii$, is monotone-carried, and $\Pii$ survives
each arrow $\Lto$ (that is, $\Pii_i \rightsquigarrow \Pii_{i+1}$ at every stage
and $\Pii_n = \Pii_\home$). If $\Mu(R_i, \text{you}) \ne 0$ for every $i$ along
the chain, then $\Mu(\home, \text{you}) \ne 0$.
\end{lemma}

\begin{proof}
Because $\Mu = \Mubar \circ \Pii$, we may read the hypothesis as
$\Mubar(\Pii_i, \text{you}) \ne 0$ for all $i$. Induct along the chain. For the
base, $\Mubar(\Pii_0, \text{you}) \ne 0$. For the step, assume
$\Mubar(\Pii_i, \text{you}) \ne 0$. The arrow $\Lto$ collapses only the
extensional state: by Proposition~\ref{prop:prov} the provenance is untouched,
so the step is the append $\Pii_i \subseteq \Pii_{i+1}$. Monotone-carry gives
$\Mubar(\Pii_{i+1}, \text{you}) \ne 0$. At the last stage
$\Pii_n = \Pii_\home$, so $\Mu(\home, \text{you}) = \Mubar(\Pii_\home,
\text{you}) \ne 0$.
\end{proof}

\begin{remark}
The proof is the whole content: $\Mu$ is constant on the fibres of $\Pii$ rather
than those of $\Rho$; $\Pii$ is the sufficient fibre; and $\Pii$ is carried, so
the coupling rides it to the limit. Nothing is claimed about $\Rho$ except that
it is the wrong fibre to read the coupling from.
\end{remark}

\section{The residue relocates}

The collapse $\Rho \downarrow 0$ is the lossy fibre. It does not destroy the
residue; it relocates it into the provenance, which is why the coupling reaches
$\home$ intact. This is the chain-level counterpart of the residue of
Chapter 1: the proof extracted from a run is a function of $\Pii$, the
extensional state is a function of the run, and neither determines the other.
A prover that collapses its state to $\Rho$ after every iteration, for memory,
can continue indefinitely and can still report that $\bot$ was derived, but
cannot say how. What survives is exactly the carried fibre, and the observer
coupling is written there.

The formal point, stated once: admissibility of a collapse is relative to a
declared task and a declared continuation; the extensional state is minimal for
all futures (Thm.~13.8); the provenance is the sufficient fibre for proof
extraction (Prop.~13.11); and an observer invariant that factors through the
carried fibre is preserved by composition to the fixed point.

\section*{Status and scope}

The lemma is proved modulo two hypotheses stated explicitly above: that $\Mu$
factors through $\Pii$, and that $\Mu$ is monotone-carried. The first is the
content of Section 2, where the factorization criterion (Prop.~1.12) and the
failure of the extensional state (Prop.~13.11) together force $\Mu$ off $\Rho$;
the second makes precise the phrase ``$\Pii$ propagates, append-only, never
lost.'' The reduction $\Rho \downarrow 0$ is the lossy fibre and the residue
relocates into $\Pii$; everything else is machinery already seated in
\emph{The Residue of Proof}. The lemma is offered as a reachable-history note,
not as a new theorem of the book.

\begin{thebibliography}{9}
\bibitem{residue}
Flyxion. \emph{The Residue of Proof: Sufficiency, History, and Repair in
Automated Reasoning}. Memnet, \texttt{automated-proofs/}, 2026.
Cross-references: Def.~13.3, Thm.~13.8, Prop.~13.11 (Chapter 13, \emph{Search
States as Worlds}); Prop.~1.12 (Chapter 1, \emph{Proof Systems and Search});
Thm.~10.12 (Chapter 10, \emph{F-Sufficiency of Certificates}).
\end{thebibliography}

\end{document}
